\documentclass[a4paper,12pt]{article}
\usepackage[utf8]{inputenc}
\usepackage[T1]{fontenc}
\usepackage[french]{babel}
\usepackage{geometry}
\geometry{left=2.5cm,right=2.5cm,top=2.5cm,bottom=2.5cm}

\usepackage{lmodern}
\usepackage{xcolor}
\definecolor{titleblue}{RGB}{0,51,140}

\usepackage{hyperref}

\pagestyle{empty}

\begin{document}

\begin{flushleft}
    \textbf{Ismail AISSAOUI} \\
    Ingénieur d'État en Intelligence Artificielle \\
    ismail.aissaoui.pro@gmail.com \\
    +213 660 70 77 96 \\
    \href{https://ismail-aissaoui.vercel.app}{ismail-aissaoui.vercel.app}
\end{flushleft}

\begin{flushright}
    À l'attention de l'équipe d'encadrement du \\
    \textbf{Laboratoire CESI LINEACT} \\
    Campus de Strasbourg \\
    \medskip
    Le 30 Avril 2026
\end{flushright}

\bigskip

\noindent \textbf{Objet : Candidature pour le doctorat en Federated GNN pour la cybersécurité des véhicules électriques}

\bigskip

Madame, Monsieur,

Ingénieur d'État en Intelligence Artificielle diplômé de l'ESI-SBA, c'est avec un vif enthousiasme que je vous adresse ma candidature pour le doctorat intitulé "Federated Graph Neural Networks for Intrusion Detection in electric vehicle networks". Passionné par la sécurité des systèmes cyber-physiques et la modélisation structurelle par graphes, ce projet au sein du consortium eTruckCharge réunit tous mes centres d'intérêt scientifiques.

Mon parcours académique et mes projets personnels m'ont permis de forger une expertise solide dans les domaines clés de cette thèse :

Premièrement, mon projet **Career Connect** a été l'occasion de concevoir une architecture de \textbf{Graph Neural Networks (GNN)} pour traiter des données relationnelles hétérogènes. J'y ai exploré l'utilisation des mesures de centralité pour améliorer la représentation des nœuds, une approche qui fait écho direct à l'originalité de vos travaux de recherche (notamment le framework GDLC mentionné dans votre offre). Je maîtrise les outils comme PyTorch Geometric et NetworkX, essentiels pour modéliser les topologies de communication véhiculaire comme des systèmes structurés en graphes.

Deuxièmement, mon expérience actuelle à la Sonatrach sur les \textbf{Jumeaux Numériques} m'a familiarisé avec la sécurisation et le monitoring en temps réel d'infrastructures critiques. En travaillant sur des modèles comme Mamba-SSM et xLSTM sous contraintes de latence (0.04ms), j'ai acquis une rigueur méthodologique cruciale pour le déploiement de modèles de détection d'intrusions sur des nœuds edge ou des chargeurs de véhicules électriques.

Enfin, je suis particulièrement motivé par la dimension collaborative et distribuée de l'apprentissage fédéré. Je suis convaincu que mon autonomie, ma capacité à assimiler rapidement des états de l'art complexes et mon goût pour la recherche appliquée me permettront de contribuer significativement aux objectifs du projet FAIR'nRG.

Intégrer le laboratoire CESI LINEACT et collaborer avec des partenaires comme ChargeMap serait pour moi un honneur et une étape majeure dans ma carrière de chercheur. Je vous remercie de l'attention que vous porterez à mon dossier.

Je vous prie d'agréer, Madame, Monsieur, l'expression de mes salutations distinguées.

\bigskip

\begin{flushright}
    \textbf{Ismail AISSAOUI}
\end{flushright}

\end{document}
